\section{Ethics and Broader Impact}

GEC explanations may influence how learners understand grammar and how writers decide whether to accept a correction. A fluent but unfaithful explanation can therefore cause educational harm even when the corrected sentence is acceptable. This work aims to reduce that risk by separating edit correspondence, model-behavior consistency, grammatical validity, and helpfulness rather than reporting a single undifferentiated score.

The pilot uses public datasets and public open-source model checkpoints with license notes recorded in the data statement. Some resources are non-commercial, so downstream use must respect those restrictions. No private user data, paid API output, or fabricated human annotation is used. The benchmark includes automatic negative examples and generated explanations; these should be used for diagnostic research and not presented to learners as authoritative feedback.
